\uuid{FEcw}
\exo7id{1599}
\titre{exo7 1599}
\auteur{ortiz}
\organisation{exo7}
\datecreate{1999-04-01}
\isIndication{false}
\isCorrection{false}
\chapitre{Réduction d'endomorphisme, polynôme annulateur}
\sousChapitre{Valeur propre, vecteur propre}
\cpge{mpsi-mp2i.applications-lineaires.b, pcsi.applications-lineaires.b, ptsi.applications-lineaires.b, mp-mpi.reduction-endomorphismes.b, pc.reduction.a, psi.reduction-endomorphismes.a, pt.reduction-endomorphismes.a, bcpst2.valeurs-propres-vecteurs-propres.a}
\module{Algèbre}
\niveau{L2}
\difficulte{}

\contenu{
\question{
Soient $E$ un K-espace vectoriel de dimension
finie et $f,g\in \mathcal{L}(E)$. Montrer que si
$\lambda$ est valeur propre de $g\circ f$ alors
$\lambda$ est valeur propre de $f\circ g$ (on
distinguera les cas $\lambda=0$ et $\lambda\not=
0$).
}
}
